\documentclass[a4paper,12pt]{article}
\usepackage[left=1.5cm, right=1.5cm, top=2cm,bottom=2cm]{geometry}
\usepackage{amsmath, amssymb, amsfonts}
\usepackage{multicol}

\usepackage{xcolor}
\usepackage{paralist}

\usepackage{enumerate}

% https://tex.stackexchange.com/questions/29834/closed-square-root-symbol
\usepackage{letltxmacro}
\makeatletter
\let\oldr@@t\r@@t
\def\r@@t#1#2{%
\setbox0=\hbox{$\oldr@@t#1{#2\,}$}\dimen0=\ht0
\advance\dimen0-0.2\ht0
\setbox2=\hbox{\vrule height\ht0 depth -\dimen0}%
{\box0\lower0.4pt\box2}}
\LetLtxMacro{\oldsqrt}{\sqrt}
\renewcommand*{\sqrt}[2][\ ]{\oldsqrt[#1]{#2}}
\makeatother

\begin{document} 
  
\section*{Piape Matemática} 
\textbf{Oficina de Operações com Gráficos}\\ 
\begin{multicols}{2}
\paragraph*{1.} Descreva verbalmente as transformações aplicadas ao gráfico de $f(x)$ para obter o gráfico de $g(x)$ em cada um dos itens abaixo:

\begin{enumerate}[a)]
    \item $g(x) = f(x) + 5$
    \item $g(x) = f(x - 3)$
    \item $g(x) = 2f(x)$
    \item $g(x) = f(4x)$
    \item $g(x) = -f(x)$
    \item $g(x) = f(-x)$
    \item $g(x) = f(x) - 10$
    \item $g(x) = \frac{1}{2}f(x)$
    \item $g(x) = f\left(\frac{1}{3}x\right)$
    \item $g(x) = -f(x + 1) + 2$
\end{enumerate} \newcolumn

\paragraph*{2.} Escreva a equação da função $g(x)$ que representa a transformação descrita em cada item, a partir da função $f(x)$:
    \begin{enumerate}[a)]
        \item Deslocar o gráfico de $f(x)$ 4 unidades para baixo.
        \item Deslocar o gráfico de $f(x)$ 2.5 unidades para a direita.
        \item Alongar verticalmente o gráfico de $f(x)$ por um fator de 3.
        \item Comprimir horizontalmente o gráfico de $f(x)$ por um fator de 2.
        \item Refletir o gráfico de $f(x)$ em torno do eixo $x$.
        \item Refletir o gráfico de $f(x)$ em torno do eixo $y$.
        \item Comprimir verticalmente o gráfico de $f(x)$ por um fator de $\frac{1}{4}$.
        \item Alongar horizontalmente o gráfico de $f(x)$ por um fator de 5.
        \item Deslocar o gráfico de $f(x)$ 7 unidades para a esquerda e 1 unidade para cima.
        \item Alongar verticalmente o gráfico de $f(x)$ por um fator de 2 e refletir em torno do eixo $y$.
    \end{enumerate}

\vspace*{0cm}
\end{multicols}
 
\vspace*{\fill}
{\footnotesize\color{gray}
\paragraph*{Gabarito} \hspace*{\fill}\\
\textbf{1.}
  \begin{inparaenum}[a)]
        \item Deslocamento vertical de 5 unidades para cima.
        \item Deslocamento horizontal de 3 unidades para a direita.
        \item Alongamento vertical por um fator de 2.
        \item Compressão horizontal por um fator de 4.
        \item Reflexão vertical (em torno do eixo $x$).
        \item Reflexão horizontal (em torno do eixo $y$).
        \item Deslocamento vertical de 10 unidades para baixo.
        \item Compressão vertical por um fator de $\frac{1}{2}$.
        \item Alongamento horizontal por um fator de 3.
        \item Reflexão vertical (em torno do eixo $x$), deslocamento horizontal de 1 unidade para a esquerda e deslocamento vertical de 2 unidades para cima.
    \end{inparaenum}
\textbf{2.}
\begin{inparaenum}[a)]
        \item $g(x) = f(x) - 4$
        \item $g(x) = f(x - 2.5)$
        \item $g(x) = 3f(x)$
        \item $g(x) = f(2x)$
        \item $g(x) = -f(x)$
        \item $g(x) = f(-x)$
        \item $g(x) = \frac{1}{4}f(x)$
        \item $g(x) = f\left(\frac{1}{5}x\right)$
        \item $g(x) = f(x + 7) + 1$
        \item $g(x) = 2f(-x)$
    \end{inparaenum}
}
\end{document}